\documentclass[11pt]{article}
\usepackage[T1]{fontenc}
\usepackage{lmodern}
\usepackage[margin=1in]{geometry}
\usepackage{amsmath,amssymb,amsthm,mathtools}
\usepackage{array,booktabs,microtype,enumitem}
\usepackage[hidelinks]{hyperref}
\numberwithin{equation}{section}
\newtheorem{theorem}{Theorem}[section]
\newtheorem{proposition}[theorem]{Proposition}
\newtheorem{lemma}[theorem]{Lemma}
\newtheorem{corollary}[theorem]{Corollary}
\theoremstyle{definition}
\newtheorem{definition}[theorem]{Definition}
\theoremstyle{remark}
\newtheorem{remark}[theorem]{Remark}
\newcommand{\HH}{\mathcal H}
\newcommand{\ceil}[1]{\left\lceil #1\right\rceil}
\newcommand{\floor}[1]{\left\lfloor #1\right\rfloor}
\newcommand{\Gap}{\mathrm{Gap}}
\allowdisplaybreaks[2]
\setlength{\emergencystretch}{2em}

\title{A six-sevenths bound for transversals of hypergraphs\\ with property $(7,2)$}
\author{}
\date{}

\begin{document}
\maketitle

\begin{abstract}
A hypergraph has property $(7,2)$ if every subfamily of at most seven
edges has a transversal of at most two points. We prove that every
finite $k$-uniform hypergraph with property $(7,2)$ has transversal
number at most $\lceil 6k/7\rceil$, for every $k\ge2$. This improves the
asymptotic coefficient $7/8$ of Fon-Der-Flaass, Kostochka and Woodall
to $6/7$ and includes the rank $k=7$, not covered by their bound; the
conjectured asymptotic value is $3k/4$. The proof refines their
avoidance method. Its main new tool is a static allocation of four
avoidance requests which, together with a symmetric allocation, closes
every ``good triple'' of edges whose pairwise intersections have at most
$T/2$ points, except those with exactly two intersections of medium
size; here $T=\lceil6k/7\rceil$. Gap information about pairwise
intersection sizes, bootstrapped in three stages, handles the remaining
triples. We also explain why $6/7$ appears to be the natural
limit of these constructions: three independent constraints of the
argument coincide exactly there.
\end{abstract}

\section{Introduction}\label{sec:intro}

A \emph{transversal} of a hypergraph is a set of points meeting every
edge, and the \emph{transversal number} $\tau(\HH)$ is the least size
of a transversal. Following Erd\H{o}s, Hajnal and Tuza~\cite{EHT}, a
hypergraph has \emph{property $(p,t)$} if every subfamily of at most $p$
edges has a transversal of size at most $t$. For integers $k\ge2$ and
$p\ge3$ let $f(k,p)$ denote the largest transversal number of a finite
$k$-uniform hypergraph with property $(p,2)$. (Using ``at most $p$
edges'' avoids vacuous exceptions for families with fewer than $p$
edges; for $p=7$ this is also the convention of the proof
in~\cite{FKW}.)

Erd\H{o}s, Fon-Der-Flaass, Kostochka and Tuza~\cite{EFKT} determined
$f(k,p)$ for $3\le p\le6$; in particular $f(k,6)=k$. For $p=7$ the
behaviour changes. Fon-Der-Flaass, Kostochka and Woodall~\cite{FKW}
proved
\[
 f(k,7)\le\ceil{7k/8}\qquad(k\ge8),
\]
and constructed $4m$-uniform families with property $(7,2)$ and
transversal number $3m+1$ (for $m\ge10$, and by their remark for $m\ge4$).
Whether $f(k,7)=(3/4+o(1))k$ is Problem~644 in Bloom's list of Erd\H{o}s
problems~\cite{Bloom}; it remains open. Kostochka~\cite{Kostochka} and
Buci\'c, Kor\'andi and Sudakov~\cite{BKS} studied the general parameters
for large $p$; their results do not apply at $p=7$.

We improve the coefficient $7/8$ to $6/7$.

\begin{theorem}\label{thm:main}
Let $k\ge2$. Every finite $k$-uniform hypergraph $\HH$ with property
$(7,2)$ satisfies
\[
 \tau(\HH)\le\ceil{6k/7}.
\]
The same bound holds for finite families of nonempty sets of size at
most $k$. In particular $f(k,7)\le\ceil{6k/7}$ for all $k\ge2$, and
$\limsup_{k\to\infty}f(k,7)/k\le6/7$.
\end{theorem}

For $2\le k\le6$ we have $\ceil{6k/7}=k$, and the theorem follows from
the upper bound $f(k,6)\le k$ of~\cite{EFKT} (see also~\cite[Section~1]{FKW}):
property $(7,2)$ implies property $(6,2)$, and this holds in either
convention, since a family with at most six edges has a transversal of
at most $2\le k$ points by property $(7,2)$ itself. The content of
Theorem~\ref{thm:main} is therefore the range $k\ge7$. The rank $k=7$
needs one additional construction (Proposition~\ref{prop:rankseven}); for
$k\ge8$ the general argument suffices. The bound $\ceil{6k/7}$ is smaller
than $\ceil{7k/8}$ for every $k\ge49$ and for many smaller $k$; the two
coincide for $21$ values of $k$ between $8$ and $48$. For $k=7$ the
previously available bound was $f(7,7)\le f(7,6)=7$.

\subsection*{The method}
The proof refines the avoidance method of~\cite{FKW}. If
$\tau(\HH)>T$, then every set of at most $T$ points is disjoint from
some edge; we call such a set a \emph{request} and such an edge a
\emph{response}. Three edges with empty common intersection form a
\emph{good triple}; its three pairwise intersections are its \emph{pair
cells}. Four further requests, chosen so that every two-point
transversal of the edges obtained so far lies inside some request,
produce a subfamily of at most seven edges without a two-point
transversal, contradicting property $(7,2)$. We say that the triple
\emph{closes}. Fon-Der-Flaass, Kostochka and Woodall used exactly this
scheme, together with case distinctions on pairwise intersection sizes
and a final argument based on a largest ``small'' intersection. We
inherit that framework. What is new is the set of closing constructions
and the way they are combined.

Put $T=\ceil{6k/7}$ and $h=k-T=\floor{k/7}$. Call a pair cell
\emph{small} if its size is at most $\gamma=T-\ceil{k/2}$, and
\emph{medium} if its size lies in $(\gamma,T/2]$. Normalized by $k$, the
thresholds are $5/14$ and $3/7$. The key new ingredient is the
following static observation (Corollary~\ref{cor:hall}): \emph{a good
triple closes by four requests fixed in advance, with no adaptivity,
whenever one pair cell has size at most $T/2$, the other two are small,
and those two have total size at most $3T-2k$.} Together with a
symmetric static allocation, this settles every triple whose cells are
at most $T/2$, except those with exactly two medium cells
(Section~\ref{sec:finish}). The two static allocations do not close the
triples with two medium cells, nor the triples containing a cell in
$(T/2,k/2]$; we treat those with adaptive requests and with information
about which intersection sizes occur. That information is produced in
three stages. The first stage uses no prior
information and excludes intersection sizes just above $T/2$. The
second uses the first to exclude an interval of small sizes. The third
uses the second to exclude the rest of $(T/2,k/2]$
(Section~\ref{sec:gaps}). The second stage and the finishing argument
have the same structure: static closure unless two cells are medium,
and a gap lemma in that case.

The coefficient $6/7$ is not an artefact of rounding. Three independent
constraints of the argument meet exactly at $T=6k/7$: the load
inequality $3k+a\le4T$ of the static allocation at $a=T/2$; the
requirement that a request built from a pair meeting in $q$ points
forces small traces, which needs $q\ge3k-3T$; and the budget
$2k+M+z\le3T$ of the near-core construction at $M=T/2$, $z=k-T$. We
explain this in Section~\ref{sec:limits}. It suggests that improving
the coefficient requires a new type of construction rather than a
retuning of parameters.

\subsection*{Organization}
Section~\ref{sec:framework} sets up requests, labels, integral
allocations and good triples. Section~\ref{sec:static} proves the two
static closures. Section~\ref{sec:adaptive} states the adaptive
closures; their proofs are in Appendix~\ref{app:adaptive}.
Section~\ref{sec:finish} proves that it suffices to exclude pair
intersections in $(T/2,k/2]$, and Section~\ref{sec:gaps} excludes them.
Section~\ref{sec:proofmain} completes the proof, and
Section~\ref{sec:limits} discusses the limits of the method. Every step is a hand argument with exact integer quantities;
no computation is used in the proof.
Table~\ref{tab:dependencies} lists, for each step of the global
argument, the lemma that closes each of its cases.

\begin{table}[ht]
\centering\small
\begin{tabular}{>{\raggedright\arraybackslash}p{0.36\textwidth}>{\raggedright\arraybackslash}p{0.54\textwidth}}
\toprule
Step & Lemma used in each case\\
\midrule
Local finisher, Proposition~\ref{prop:localfinish}
 & Corollary~\ref{cor:hall}; Lemma~\ref{lem:symmetric}; Lemma~\ref{lem:twocores}; Lemma~\ref{lem:nearcore}\\
Small maximum, Theorem~\ref{thm:finish}
 & Corollary~\ref{cor:hall} (twice); Lemma~\ref{lem:twolarge}\\
Rank seven, Proposition~\ref{prop:rankseven}
 & explicit requests; Corollary~\ref{cor:hall}\\
First gap, Proposition~\ref{prop:gapone}
 & Lemma~\ref{lem:splitA}; Lemma~\ref{lem:onecore}; Lemma~\ref{lem:symmetric}\\
Middle gap, Proposition~\ref{prop:gaptwo}
 & Corollary~\ref{cor:hall}; Lemma~\ref{lem:twocores}\\
Upper gap, Proposition~\ref{prop:gapthree}
 & Lemma~\ref{lem:splitB}; Lemma~\ref{lem:gaptraces}; Lemma~\ref{lem:gapcore}\\
\bottomrule
\end{tabular}
\caption{Case structure of the proof. Each case is closed by one lemma.}
\label{tab:dependencies}
\end{table}

\section{The framework}\label{sec:framework}

Throughout Sections~\ref{sec:framework}--\ref{sec:adaptive}, $\HH$ is a
finite $k$-uniform hypergraph and $T$ is an integer with $0\le T\le k$
and $\tau(\HH)>T$. A \emph{request} is a set of at most $T$ points. Since
$\tau(\HH)>T$, every request is disjoint from some edge of $\HH$, a
\emph{response} to it. A subfamily is \emph{bad} if it has no
transversal of size at most two. A bad subfamily of at most seven edges
contradicts property~$(7,2)$.

\subsection{Covering by requests}
\begin{lemma}[Request-cover principle]\label{lem:requestcover}
Let $A_1,\dots,A_t$ be edges and $D_1,\dots,D_r$ requests, $t+r\le7$.
If every transversal of $\{A_1,\dots,A_t\}$ of size at most two is
contained in some $D_j$, then $A_1,\dots,A_t$ together with responses
to $D_1,\dots,D_r$ form a bad subfamily of at most seven edges.
\end{lemma}
\begin{proof}
A transversal of size at most two of the whole subfamily meets
$A_1,\dots,A_t$, so it lies in some $D_j$ and misses the response to
$D_j$. Repeated edges only decrease the number of distinct edges.
\end{proof}

Requests may be chosen adaptively: a request may depend on earlier
responses. A transversal of size one is contained in a transversal of
size two whenever there are at least two points, so we only discuss
pairs of points. A pair $\{u,v\}$ of points is a \emph{candidate pair}
of a family of edges if it meets every edge of the family.

\emph{Labels.} Given requests $R_1,\dots,R_r$, the \emph{label} of a
point is the set of indices $j$ with the point in $R_j$. A pair $\{u,v\}$
lies in a common request exactly when the labels of $u$ and $v$
intersect. We often partition the points into \emph{cells} and
prescribe the labels allowed in each cell; then two cells whose points
form candidate pairs must receive labels that pairwise intersect.

\subsection{Integral allocation}
Most constructions place some cells in prescribed requests and then
distribute other cells among several allowed requests. The following
form of the supply--demand theorem guarantees integral distributions.

\begin{lemma}[Integral allocation]\label{lem:allocation}
Let $W_1,\dots,W_n$ be disjoint finite sets. Suppose each $W_i$ has a
set $N_i\subseteq\{1,\dots,r\}$ of allowed indices, and each index $j$
has an integer capacity $c_j\ge0$. The points of $\bigcup W_i$ can be
assigned so that every point of $W_i$ receives one index of $N_i$ and
index $j$ receives at most $c_j$ points if and only if
\[
 \sum_{i\in I}|W_i|\le\sum_{j\in N(I)}c_j
 \qquad\text{for every } I\subseteq\{1,\dots,n\},\quad
 N(I)=\bigcup_{i\in I}N_i .
\]
\end{lemma}
\begin{proof}
Necessity is clear. For sufficiency, form the network with arcs from a
source to $W_i$ of capacity $|W_i|$, arcs of infinite capacity from $W_i$
to each $j\in N_i$, and arcs from $j$ to a sink of capacity $c_j$.
Consider a cut of finite capacity and let $I$ be the set of indices $i$
such that $W_i$ lies on the source side. No arc of infinite capacity is
cut, so $N(I)$ lies on the source side as well, and the capacity of the
cut is at least $\sum_{i\notin I}|W_i|+\sum_{j\in N(I)}c_j$, which by
hypothesis is at least $\sum_i|W_i|$. By the max-flow min-cut theorem and the integrality
of maximum flows with integer capacities, there is an integral flow
saturating every source arc; it gives the assignment.
\end{proof}

When $n=1$ the condition is $|W_1|\le\sum_{j\in N_1}c_j$. Adding a
point to a request never destroys a covering, so a construction may
always put a cell into more requests than necessary, as long as the
sizes stay at most $T$.

\subsection{Good triples}
A triple $E,F,G$ of edges is \emph{good} if $E\cap F\cap G=\varnothing$.
Throughout we write
\[
 X=E\cap F,\quad Y=E\cap G,\quad Z=F\cap G,\qquad
 x=|X|,\ y=|Y|,\ z=|Z|,\ S=x+y+z,
\]
and $P_E=E\setminus(F\cup G)$, $P_F=F\setminus(E\cup G)$,
$P_G=G\setminus(E\cup F)$, of sizes $k-x-y$, $k-x-z$, $k-y-z$. Because
the triple is good, the six cells $X,Y,Z,P_E,P_F,P_G$ are disjoint.
A pair meets all three edges exactly when it lies in two different pair
cells, or in a pair cell and the opposite private part; so the candidate
pairs are
\begin{equation}\label{eq:triplecand}
 X\times Y,\ X\times Z,\ Y\times Z,\qquad
 X\times P_G,\ Y\times P_F,\ Z\times P_E .
\end{equation}
Points outside $E\cup F\cup G$ belong to no candidate pair.

A triple \emph{closes at budget $T$} if four further requests of size at
most $T$, chosen possibly adaptively, force a bad subfamily of at most
seven edges. Since only sizes of cells enter the constructions, the
statement that a triple with given pair sizes closes does not depend on
the names of its edges; we permute $E,F,G$ freely and say which cell
plays which role.

After a fourth edge $K$ is obtained, candidate pairs of $\{E,F,G,K\}$
are described as follows. Label each point by the set of edges among
$E,F,G,K$ containing it.

\begin{lemma}[Four edges]\label{lem:foursets}
Let $E,F,G$ be a good triple and $K$ an edge. The candidate pairs of
$\{E,F,G,K\}$ are exactly the pairs in
\begin{gather*}
 (X\cap K)\times G,\qquad (Y\cap K)\times F,\qquad (Z\cap K)\times E,\\
 (X\setminus K)\times(K\cap P_G),\quad
 (Y\setminus K)\times(K\cap P_F),\quad
 (Z\setminus K)\times(K\cap P_E).
\end{gather*}
\end{lemma}
\begin{proof}
A pair meets the four edges exactly when the union of the labels of its
points is $\{E,F,G,K\}$. No label contains $\{E,F,G\}$. If one label has
three elements, it is $\{E,F,K\}$, $\{E,G,K\}$ or $\{F,G,K\}$ and the
other point must lie in the missing edge; this gives the first row. If
both labels have at most two elements, they are complementary
two-element sets, and one of them contains $K$; this gives the second
row. A label with at most one element can only be completed by a
three-element label.
\end{proof}

The sets $X\cap K$, $Y\cap K$, $Z\cap K$ are the \emph{triple cells} of
the four edges. A response might coincide with an earlier edge (for
instance when its traces on the earlier edges are not forced to be smaller
than $k$). This is harmless: Lemma~\ref{lem:requestcover} counts distinct
edges, and Lemma~\ref{lem:foursets} holds also when $K$ is one of $E,F,G$. Most adaptive constructions first request an edge $K$
avoiding all or most of the pair cells, so that few triple cells
survive.

\subsection{Intersection gaps}
For integers $0\le\ell\le u<k$ we write $\Gap(\ell,u)$ for the statement
that no two distinct edges of $\HH$ meet in a number of points in
$\{\ell+1,\dots,u\}$; equivalently, $|A\cap B|\le u$ implies
$|A\cap B|\le\ell$ for distinct edges $A,B$. Because $u<k$, an edge
meeting a given edge in at most $u$ points is distinct from it, so
$\Gap(\ell,u)$ applies to every trace of a response that we bound by
$u$. We also use the \emph{dichotomy with threshold $m$}: every two
distinct edges meet in at most $m$ or in more than $k/2$ points. Two
disjoint subsets of one edge cannot both have more than $k/2$ points.

\subsection{Seed requests}
\begin{lemma}[Balanced request]\label{lem:balanced}
If two edges $E,F$ meet in $M\le T$ points, a response $G$ to a suitable
request of size $T$ forms a good triple with them and satisfies
\[
 |E\cap G|,\ |F\cap G|\le k-\floor{(T+M)/2}.
\]
\end{lemma}
\begin{proof}
Request $E\cap F$ together with $\floor{(T-M)/2}$ points of $E\setminus F$
and $\ceil{(T-M)/2}$ points of $F\setminus E$; these fit since $T\le k$.
The response avoids $E\cap F$, and each of its traces is at most
$k-M-\floor{(T-M)/2}=k-\floor{(T+M)/2}$.
\end{proof}

\begin{lemma}[Capped request]\label{lem:caps}
Suppose two edges meet in $q$ points and $K\ge0$ is an integer with
$0\le2k-T-q\le2K$ and $K\le k-q$. Then a response to a suitable request
of size $T$ forms a good triple with the two edges, with both new pair
intersections at most $K$.
\end{lemma}
\begin{proof}
Choose integers $u,v\in[0,K]$ with $u+v=2k-T-q$. Keep $u$ points of the
first private part and $v$ of the second (possible as $K\le k-q$) and
request all other points of the union of the two edges, which has
$2k-q-u-v=T$ points and contains their intersection.
\end{proof}

\section{Static closures}\label{sec:static}

A static closure uses four requests fixed in advance. By
Lemma~\ref{lem:requestcover} and~\eqref{eq:triplecand}, four requests
close a good triple when the three pair cells receive pairwise
intersecting labels and each private part receives labels meeting every
label of the opposite pair cell.

\begin{lemma}[Hall allocation]\label{lem:hall}
Let a good triple have pair sizes $a=|X|$, $b=|Y|$, $c=|Z|$. It closes at
budget $T$ by a static allocation if
\begin{gather}
 T-a\ge0,\quad T-a-b\ge0,\quad T-a-c\ge0,\quad T-b-c\ge0,\label{eq:hallbase}\\
 k+2b\le2T,\quad k+2c\le2T,\quad k+3a\le3T,\quad 2k+b+c\le3T,\label{eq:hallseven}\\
 2k+2a+b\le4T,\quad 2k+2a+c\le4T,\quad 3k+a\le4T.\notag
\end{gather}
\end{lemma}
\begin{proof}
Give the points of $X$, $Y$, $Z$ the labels $\{1,2,3\}$, $\{2,4\}$,
$\{3,4\}$; these pairwise intersect. Every point of $P_G$ will receive
one of $\{1\},\{2\},\{3\}$, every point of $P_F$ one of $\{2\},\{4\}$,
and every point of $P_E$ one of $\{3\},\{4\}$. Each of these meets the
label of the opposite pair cell, so all candidate
pairs~\eqref{eq:triplecand} are covered. After the pair cells are placed,
the remaining capacities of the four requests are $T-a$, $T-a-b$,
$T-a-c$, $T-b-c$, which are nonnegative by~\eqref{eq:hallbase}. The
three private parts have sizes $k-a-b$, $k-a-c$, $k-b-c$ and allowed
index sets $\{3,4\}$, $\{2,4\}$, $\{1,2,3\}$. The seven conditions of
Lemma~\ref{lem:allocation} are, for $P_E$, $P_F$, $P_G$, $P_E\cup P_F$,
$P_E\cup P_G$, $P_F\cup P_G$ and all three, exactly the seven
inequalities~\eqref{eq:hallseven}; for example, for $P_G$ the condition
is $k-b-c\le(T-a)+(T-a-b)+(T-a-c)$, that is $k+3a\le3T$, and for all
three it is $3k-2a-2b-2c\le4T-3a-2b-2c$, that is $3k+a\le4T$.
\end{proof}

The next corollary is the form in which the Hall allocation is used.

\begin{corollary}\label{cor:hall}
Let $T\ge6k/7$. A good triple closes at budget $T$ if one pair cell has
size $a\le T/2$ and the other two have sizes $b,c\le T-k/2$ with
$b+c\le3T-2k$.
\end{corollary}
\begin{proof}
We check Lemma~\ref{lem:hall}. We have $T-a\ge T/2$,
$T-a-b\ge T-T/2-(T-k/2)=(k-T)/2\ge0$, likewise $T-a-c\ge0$, and
$T-b-c\ge T-(3T-2k)=2(k-T)\ge0$. Next $k+2b,\,k+2c\le2T$ and $2k+b+c\le3T$
hold by hypothesis. Finally, using $T\ge6k/7$,
\begin{gather*}
 k+3a\le k+\tfrac32T\le3T,\qquad
 2k+2a+b\le2k+T+(T-\tfrac k2)=\tfrac32k+2T\le4T,\\
 3k+a\le3k+\tfrac T2\le4T .
\end{gather*}
The last inequality is equivalent to $7T\ge6k$.
\end{proof}

The second static closure treats triples whose three cells are all of
roughly the same size. It rests on an exact allocation lemma.

\begin{lemma}[Boxed triangle]\label{lem:triangle}
Let $M_1,M_2,M_3\ge0$ and $L_1,L_2,L_3$ be integers with
$L_i\le M_j+M_l$ whenever $\{i,j,l\}=\{1,2,3\}$, and let $t\ge0$ be an
integer. There are integers $0\le a_i\le M_i$ with $a_j+a_l\ge L_i$ for
all $\{i,j,l\}=\{1,2,3\}$ and $a_1+a_2+a_3\le t$ if and only if
\[
 t\ge\max\Bigl\{L_1,L_2,L_3,\ L_1+L_2-M_3,\ L_1+L_3-M_2,\ L_2+L_3-M_1,\
 \tfrac12(L_1+L_2+L_3)\Bigr\}.
\]
\end{lemma}
\begin{proof}
If such $a_i$ exist with $\sigma=a_1+a_2+a_3\le t$, then $\sigma\ge L_i$
because $a_i\ge0$; adding $a_2+a_3\ge L_1$ and $a_1+a_3\ge L_2$ gives
$\sigma+a_3\ge L_1+L_2$, so $\sigma\ge L_1+L_2-M_3$, and similarly for the
other two pairs; and adding all three constraints gives
$2\sigma\ge L_1+L_2+L_3$.

Conversely, suppose $t$ satisfies the displayed inequality. Every
displayed quantity is at most $M_1+M_2+M_3$: for instance
$L_1+L_2-M_3\le(M_2+M_3)+(M_1+M_3)-M_3$ and
$L_1+L_2+L_3\le2(M_1+M_2+M_3)$. So $t'=\min(t,M_1+M_2+M_3)$ also
satisfies it, and we look for $a_i$ with $a_1+a_2+a_3=t'$. Then the
constraints become $0\le a_i\le\min(M_i,t'-L_i)$, and integers with sum
$t'$ exist if and only if each upper bound is nonnegative and the three
upper bounds sum to at least $t'$. The first requirement is $t'\ge L_i$.
Expanding the sum of three minima as the minimum of eight sums, the
second requirement is $M_1+M_2+M_3\ge t'$, $M_j+M_l+t'-L_i\ge t'$ (the
host conditions), $M_i+2t'-L_j-L_l\ge t'$, and $3t'-L_1-L_2-L_3\ge t'$.
All hold.
\end{proof}

\begin{lemma}[Symmetric allocation]\label{lem:symmetric}
Let a good triple have pair sizes $m\ge y\ge z$ (in any order of the
cells) and $S=m+y+z$. It closes at budget $T$ by a static allocation if
\[
 T\ge k+m-y-z,\qquad 2T\ge k+m,\qquad 3T\ge2k+m+y-z,\qquad 5T\ge3k+S.
\]
\end{lemma}
\begin{proof}
Call the pair cells $\mathsf M,\mathsf Y,\mathsf Z$ (sizes $m,y,z$) and
the opposite private parts $P_{\mathsf M},P_{\mathsf Y},P_{\mathsf Z}$.
Split $\mathsf M=\mathsf M_1\sqcup \mathsf M_2$,
$\mathsf Y=\mathsf Y_1\sqcup \mathsf Y_2$,
$\mathsf Z=\mathsf Z_1\sqcup \mathsf Z_2$ with
$|\mathsf M_1|=a_1$, $|\mathsf Y_1|=a_2$, $|\mathsf Z_1|=a_3$, and use
the requests
\begin{align*}
 R_0&=\mathsf M_1\cup \mathsf Y_1\cup \mathsf Z_1,&
 R_{\mathsf M}&=\mathsf M\cup P_{\mathsf M}\cup \mathsf Y_2\cup \mathsf Z_2,\\
 R_{\mathsf Y}&=\mathsf Y\cup P_{\mathsf Y}\cup \mathsf M_2\cup \mathsf Z_2,&
 R_{\mathsf Z}&=\mathsf Z\cup P_{\mathsf Z}\cup \mathsf M_2\cup \mathsf Y_2.
\end{align*}
The parts $\mathsf M_1,\mathsf Y_1,\mathsf Z_1$ have labels
$\{0,\mathsf M\},\{0,\mathsf Y\},\{0,\mathsf Z\}$, the parts
$\mathsf M_2,\mathsf Y_2,\mathsf Z_2$ all have label
$\{\mathsf M,\mathsf Y,\mathsf Z\}$, and $P_{\mathsf M},P_{\mathsf Y},
P_{\mathsf Z}$ have labels $\{\mathsf M\},\{\mathsf Y\},\{\mathsf Z\}$.
Any two labels of different pair cells intersect, and each private label
meets both labels of its opposite pair cell, so all candidate pairs are
covered. The sizes of the four requests are
\[
 a_1+a_2+a_3,\quad k+m-a_2-a_3,\quad k+y-a_1-a_3,\quad k+z-a_1-a_2 .
\]
Hence it suffices to find integers $0\le a_1\le m$, $0\le a_2\le y$,
$0\le a_3\le z$ with $a_2+a_3\ge k+m-T$, $a_1+a_3\ge k+y-T$,
$a_1+a_2\ge k+z-T$ and $a_1+a_2+a_3\le T$. Apply
Lemma~\ref{lem:triangle} with $(M_1,M_2,M_3)=(m,y,z)$,
$(L_1,L_2,L_3)=(k+m-T,k+y-T,k+z-T)$ and $t=T$. The host conditions
reduce to $T\ge k+m-y-z$ (the other two are weaker since $m\ge y\ge z$).
The largest single $L_i$ gives $2T\ge k+m$; the largest of the three
two-term expressions is $L_1+L_2-M_3$, which gives $3T\ge2k+m+y-z$;
and the last expression gives $5T\ge3k+S$.
\end{proof}

\section{Adaptive closures}\label{sec:adaptive}

The following lemmas each describe one adaptive construction: a first
request produces a fourth edge $K$, and three further requests,
depending on $K$, cover all candidate pairs of $\{E,F,G,K\}$ given by
Lemma~\ref{lem:foursets}. The hypotheses are exact integer conditions.
Proofs are given in Appendix~\ref{app:adaptive}. We use the notation of
Section~\ref{sec:framework} and assume $T\le k$ throughout; each lemma
may be applied to any assignment of the names $E,F,G$ to the edges of a
good triple, which permutes $x,y,z$.

\subsection{Without gap information}
The first lemma requests an edge avoiding two pair cells and as much of
the third as possible; at most one triple cell survives. The next two
avoid all three pair cells.

\begin{lemma}[One surviving cell]\label{lem:onecore}
A good triple closes at budget $T$ if
\begin{gather*}
 T\ge x+y,\quad 2T\ge k+2x,\quad 2T\ge k+2y,\quad
 T\ge k+x-y-z,\quad T\ge k-x+y-z,\\
 3T\ge3k-S,\quad 5T\ge3k+S,\quad 3T\ge k+x+y+2z,\quad 4T\ge2k+3z.
\end{gather*}
\end{lemma}

\begin{lemma}[First asymmetric split]\label{lem:splitA}
A good triple closes at budget $T$ if
\[
 T\ge S,\qquad 2T\ge k+2y,\qquad 2T\ge k+2x-y+z,\qquad 3T\ge k+2x+y+3z.
\]
\end{lemma}

\begin{lemma}[Second asymmetric split]\label{lem:splitB}
A good triple closes at budget $T$ if
\[
 T\ge S,\qquad 3T\ge k+3x,\qquad T\ge k-x+z,\qquad 3T\ge k+2x+3y+z.
\]
\end{lemma}

\subsection{With gap information}
In the next three lemmas $\Gap(\ell,u)$ holds for integers
$0\le\ell\le u<k$. The first request is designed so that some traces of
$K$ are at most $u$, and hence at most $\ell$.

\begin{lemma}[Two forced traces]\label{lem:gaptraces}
Assume $\Gap(\ell,u)$ and put $p=(k-x-y-u)_+$ and $t=(k-x-z-u)_+$.
A good triple closes at budget $T$ if
\[
 S+p+t\le T,\qquad k+3x+p+t\le3T,\qquad y+z+2\ell\le T.
\]
\end{lemma}

\begin{lemma}[Forced traces and one core]\label{lem:gapcore}
Assume $\Gap(\ell,u)$ and put $\Delta=(S-T)_+$. A good triple closes at
budget $T$ if
\begin{gather*}
 x+y\le T,\qquad k-x-y\le u,\qquad k-x-z+\Delta\le u,\\
 2x+2\Delta+k-y-z\le2T,\qquad y+z+2\ell\le T,\qquad 2k+x-y+\ell+\Delta\le3T.
\end{gather*}
\end{lemma}

\begin{lemma}[Two cores]\label{lem:twocores}
Assume $\Gap(\ell,u)$ and put $g=(k-y-u)_+$. A good triple closes at
budget $T$ if
\begin{gather*}
 g\le\min(x,z),\qquad y+2g\le T,\qquad x+\ell\le T,\qquad z+\ell\le T,\\
 k+2y\le2T,\qquad 3k\le4T,\qquad 3k+S\le5T.
\end{gather*}
\end{lemma}

The last two lemmas use the dichotomy with threshold $m$, where
$2m\le k$.

\begin{lemma}[Near core]\label{lem:nearcore}
Assume the dichotomy with threshold $m$ and put $\Delta=(S-T)_+$. A good
triple closes at budget $T$ if
\begin{gather*}
 y+z\le T,\quad m+z\le T,\quad x+m\le T,\quad k+y-z\le T,\\
 2k+x-2z\le2T,\quad 2k+x+m-z\le3T,
\end{gather*}
and either
\begin{gather*}
 k-x+y+\Delta\le T,\quad 2k-2x+z+\Delta\le2T,\quad
 2k-z+\Delta\le2T,\quad 3k-x-y+\Delta\le3T,
\end{gather*}
or $S\le T$, $k-x+y\le T$ and $2k-2x+z\le2T$.
\end{lemma}

\begin{lemma}[Two large opposite cells]\label{lem:twolarge}
Assume the dichotomy with threshold $m$, where $4m\le k$. Let $E,F,G,H$
be edges with all four triple intersections empty, such that
$X=E\cap F$ and $B=G\cap H$ have sizes $x,b>k/2$ and the other four
pairwise intersections have at most $m$ points. Three further requests
force a bad subfamily of at most seven edges if
\[
 T\ge\ceil{k/2}+2m,\qquad T\ge x+m,\qquad 3T\ge2x+b+\ceil{k/2}+2m.
\]
\end{lemma}

\section{The finishing theorem}\label{sec:finish}

In this section we show that it is enough to exclude pair intersections
of size in $(T/2,k/2]$. Write
\[
 h=k-T,\qquad \gamma=T-\ceil{k/2}.
\]
An integer $w$ satisfies $w\le\gamma$ if and only if $w\le T-k/2$. We
call a pair cell \emph{small} if its size is at most $\gamma$ and
\emph{medium} if its size lies in $(\gamma,T/2]$.

\begin{proposition}[Local finisher]\label{prop:localfinish}
Let $6k/7\le T\le k$ and let $M\le T/2$ be an integer. Assume the
dichotomy with threshold $M$. Then every good triple with pair sizes
$M\ge y\ge z$ closes at budget $T$.
\end{proposition}

\begin{proof}
We say that the triple has \emph{two medium cells} if $y>\gamma\ge z$.
(Its cell of size $M$ is then medium as well.) There are four cases.

\paragraph{Case 1: not two medium cells, and $y+z\le3T-2k$.}
If $z>\gamma$ then $y+z\ge2\gamma+2\ge2T-k+1>3T-2k$, so here $z\le y\le\gamma$.
Corollary~\ref{cor:hall} applies with $a=M$, $b=y$, $c=z$.

\paragraph{Case 2: not two medium cells, and $y+z>3T-2k$.}
Either $y\le\gamma$, or $z>\gamma$. We check
Lemma~\ref{lem:symmetric} with $m=M$. First, $M\le T/2\le4T-3k$ because
$7T\ge6k$, so $k+M-T\le3T-2k<y+z$. Next $2T\ge k+T/2\ge k+M$. For the
third and fourth conditions we show
\[
 M+y-z<k-T/2\qquad\text{and}\qquad S\le3T/2 .
\]
If $z>\gamma$, then $y-z<M-(T-k/2)$, so $M+y-z<2M-T+k/2\le k/2\le k-T/2$;
and $S\le3M\le3T/2$. If $y\le\gamma$, then $y-z=2y-(y+z)<2(T-k/2)-(3T-2k)=k-T$,
so $M+y-z<k-T/2$; and $S\le T/2+2(T-k/2)\le3T/2$. Consequently
$2k+M+y-z<3k-T/2\le3T$ and $3k+S\le3k+3T/2\le5T$, both because $7T\ge6k$.

\paragraph{Case 3: two medium cells and $z>h$.}
Apply Lemma~\ref{lem:twocores} with the cell of size $z$ in the role of
$Y$ and the cells of sizes $M,y$ in the roles of $X,Z$; that is, with
$(x,y,z)$ of the lemma equal to $(M,z,y)$, and with $\Gap(M,\floor{k/2})$,
which is implied by the dichotomy (note $0\le M\le\floor{k/2}<k$). Here
$g=(k-z-\floor{k/2})_+=\ceil{k/2}-z$, as $z\le\gamma<\ceil{k/2}$. Since
$M+z\ge y+z>(T-k/2)+(k-T)=k/2$, both $M$ and $y$ are at least
$\ceil{k/2}-z=g$. The first request has size
$z+2g=\max\{z,\,k-z+(k-2\floor{k/2})\}\le T$, because $z\ge h+1$. The
remaining conditions are
\[
 M+M\le T,\quad y+M\le T,\quad k+2z\le k+2\gamma\le2T,\quad 3k\le4T,
\]
and $3k+S\le3k+2M+\gamma\le3k+T+T-k/2=\tfrac52k+2T\le5T$.

\paragraph{Case 4: two medium cells and $z\le h$.}
Apply Lemma~\ref{lem:nearcore} with $m=M$ and $(x,y,z)$ of the lemma equal
to $(M,z,y)$. With $c=T-k/2$ we have $y>c$, $z\le h$ and $h\le c$. The
common conditions are
\begin{gather*}
 z+y\le h+T/2=k-T/2\le T,\qquad M+y\le2M\le T,\qquad M+M\le T,\\
 k+z-y<k+h-c=\tfrac52k-2T\le T,\qquad
 2k+M-2y<2k+\tfrac T2-2c=3k-\tfrac32T\le2T,\\
 2k+2M-y<2k+T-c=\tfrac52k\le3T .
\end{gather*}
If $S\le T$ we use the second alternative:
$k-M+z\le k-y+z<T$ and $2k-2M+y\le2k-M<2k-c=\tfrac52k-T\le2T$. If $S>T$
then $\Delta=S-T$ and the four remaining conditions become
\begin{align*}
 k+y+2z&\le k+\tfrac T2+2h=3k-\tfrac32T\le2T,&
 2k-M+2y+z&\le2k+y+z\le3k-\tfrac T2\le3T,\\
 2k+M+z&\le2k+\tfrac T2+h=3k-\tfrac T2\le3T,&
 3k+y&\le3k+\tfrac T2\le4T,
\end{align*}
where we used $y\le M\le T/2$ and $z\le h$; each final inequality is
equivalent to $T\ge6k/7$ or weaker.

The four cases cover all triples.
\end{proof}

\begin{theorem}[Finishing theorem]\label{thm:finish}
Let $k\ge7$ and $T=\ceil{6k/7}$. If every two distinct edges of a finite
$k$-uniform hypergraph $\HH$ with property $(7,2)$ meet in at most $T/2$
or in more than $k/2$ points, then $\tau(\HH)\le T$.
\end{theorem}

\begin{proof}
Suppose $\tau(\HH)>T$. An edge avoiding a $T$-subset of another edge
meets it in at most $h=k-T\le k/2$ points, so the largest size $M$ of a
pair intersection that is at most $k/2$ exists. By hypothesis
$M\le T/2$, and by the choice of $M$ the dichotomy with threshold $M$
holds. Write $k=7h+s$ with $0\le s\le6$; then $h=\floor{k/7}\ge1$.

\emph{Case $M\ge h+1$.} (For $k=7$ this is the case $M\in\{2,3\}$.) Apply
Lemma~\ref{lem:balanced} to a pair attaining $M$. As $T+M\ge k+1$, the new pair intersections are at most
$k-\ceil{k/2}=\floor{k/2}$, hence at most $M$. The resulting good
triple closes by Proposition~\ref{prop:localfinish}, a contradiction.

\emph{Case $M\le h$ and $k=7$.} Then $h=1$, so every two edges meet in
at most one or at least four points, and
Proposition~\ref{prop:rankseven} below gives a contradiction.

\emph{Case $M\le h$ and $k\ge8$.} Then $h+s\ge2$. We record $\gamma\ge h$, $2h\le3T-2k$ and $4h\le k$,
which follow from $k\ge7h$. Apply Lemma~\ref{lem:balanced} to edges
$E,G$ with $|E\cap G|=M$, obtaining $F$. If both $|E\cap F|$ and
$|F\cap G|$ are at most $k/2$, they are at most $M$, and the triple
closes by Corollary~\ref{cor:hall}, as all three cells are at most
$h\le\gamma$ and the two cells other than $E\cap G$ sum to at most
$2h\le3T-2k$. Otherwise exactly one of them exceeds $k/2$, since they
are disjoint subsets of $F$; by symmetry
\[
 x=|E\cap F|>k/2,\qquad |E\cap G|=M,\qquad |F\cap G|\le M,
\]
and $x\le k-\floor{(T+M)/2}\le k-(T+M-1)/2$. Then
\[
 x+2M\le k-\tfrac T2+\tfrac32M+\tfrac12\le\tfrac k2+2h+\tfrac12\le T,
\]
using $T=k-h$, $M\le h$ and $k\ge6h+1$. Request $X\cup Y\cup Z$, which has
$x+|Y|+|Z|\le x+2M\le T$ points, together with $T-x-|Y|-|Z|$ private
points of $G$ (there are enough, as $|P_G|=k-|Y|-|Z|$), and let $H$ be a
response. All triple intersections of $E,F,G,H$ are empty. Since $H$
avoids $X$, its traces on $E$ and $F$ have at most $k-x<k/2$ points,
hence at most $M$. Its trace on $G$ has $b\le k-T+x$ points.

If $b\le k/2$, then $b\le M$ and $E,G,H$ is a good triple (as $H$ avoids
$Y=E\cap G$) with all cells at most $M\le h$; it closes by
Corollary~\ref{cor:hall} exactly as above. We discard $F$: the resulting
bad subfamily consists of $E,G,H$ and four responses.

If $b>k/2$, apply Lemma~\ref{lem:twolarge} to $E,F,G,H$ with $m=M$. Its
conditions hold:
\begin{gather*}
 \ceil{k/2}+2M\le\tfrac{k+1}2+2h\le k-h=T,\qquad
 x+M\le k-\tfrac T2+\tfrac M2+\tfrac12\le\tfrac k2+h+\tfrac12\le T,\\
 2x+b+\ceil{k/2}+2M\le3x+h+\tfrac{k+1}2+2M
 \le\tfrac72k-\tfrac32T+\tfrac M2+h+2\le2k+3h+2\le3T,
\end{gather*}
where the final step is $k\ge6h+2$, that is $h+s\ge2$. The resulting bad
subfamily has seven edges. In every case property $(7,2)$ is violated.
\end{proof}

At $k=7$ the last inequality fails: for $x=|E\cap F|=4$, $b=|G\cap H|=5$
and $M=1$, the third condition of Lemma~\ref{lem:twolarge} reads
$18\ge19$. A differently aimed fourth request repairs this.

\begin{proposition}[Rank seven]\label{prop:rankseven}
Let $\HH$ be a finite $7$-uniform hypergraph with property $(7,2)$ in
which every two distinct edges meet in at most one or at least four
points. Then $\tau(\HH)\le6$.
\end{proposition}

\begin{proof}
Suppose $\tau(\HH)>6$, so requests have at most six points. By
hypothesis, a response that misses some point of an earlier edge meets
that edge in $0$, $1$, or at least $4$ points; in particular,
\emph{if it meets the edge in at most three points, then it meets it in
at most one point}. An edge avoiding six points of another edge meets it
in at most one point.

\emph{Case 1: no two edges meet in exactly one point.} Then some edges
$E_1,E_2$ are disjoint: an edge avoiding six points of another meets it
in at most one point, hence in none. Request three points of each. The response $G$
meets each $E_i$ within four points, hence in $0$ or $4$ points, and not
in $4$ points of both since $|G|=7$. If $G$ is disjoint from both, the
edges $E_1,E_2,G$ are pairwise disjoint and form a bad subfamily.
Otherwise, say, $Q=G\cap E_1$ has four points and $G\cap E_2=\varnothing$.
Every candidate pair of $\{E_1,E_2,G\}$ has a point of $E_2$, and its
other point lies in $E_1\cap G=Q$. Partition $E_2$ into sets
$S_1,\dots,S_4$ of sizes $2,2,2,1$; the four requests $Q\cup S_i$ cover
$Q\times E_2$.

\emph{Case 2: two edges $E,F$ meet in one point $p$.} Request $p$, three
points of $E\setminus F$ and two points of $F\setminus E$. The response
$G$ avoids $p$, so $E,F,G$ is a good triple; $|G\cap E|\le1$ since $G$
meets $E$ within three points, and $|G\cap F|\in\{0,1,4\}$. If
$|G\cap F|\le1$, all three pair cells have at most one point and the
triple closes by Corollary~\ref{cor:hall} ($T=6=6k/7$, and
$a\le3$, $b,c\le1\le T-k/2$, $b+c\le2\le3T-2k$).

It remains to treat $|G\cap F|=4$. Rename $A=F$, $B=G$, $C=E$. Then
$X=A\cap B$ has four points, $A\cap C=\{y\}$ with $y=p$, $Z=B\cap C$ has
at most one point, and $A\cap B\cap C=\varnothing$. Put $A'=A\setminus B$
(three points, including $y$), $B'=B\setminus A$ (three points, including
$Z$) and $P_C=C\setminus(A\cup B)$ ($5$ or $6$ points). Choose
$x_1\in X$, $c_1\in P_C$ and $b^*\in B'$ with $Z\subseteq\{b^*\}$, and
request
\[
 D=\{x_1\}\cup A'\cup\{b^*,c_1\}\qquad(|D|=6).
\]
Let $H$ be a response; it differs from $A,B,C$. Its traces satisfy:
$H\cap A\subseteq X\setminus\{x_1\}$, so $|H\cap A|\le1$ and $H\cap A=H\cap X$;
$H\cap B\subseteq(X\setminus\{x_1\})\cup(B'\setminus\{b^*\})$ has at most
$1+2$ points, so $|H\cap B|\le1$; and $H\cap C\subseteq P_C\setminus\{c_1\}$,
so $\eta=|H\cap C|\in\{0,1,4,5\}$. Consequently $H$ meets at most one of
$X$ and $B'$, each in at most one point.

A candidate pair of $\{A,B,C,H\}$ contains a point of $X$, or consists of
$u\in A'$ and $v\in B'$. In the first case the other point lies in $C$
(as $X\cap C=\varnothing$), and in $H$ unless the $X$-point is in $H$. In
the second case $u\notin H$, so $v\in H\cap B'$; then $v\ne b^*$, so
$v\notin C$, and $u\in C$ forces $u=y$. Hence the candidate pairs are
\[
 \{x\}\times C\ \text{ if } H\cap X=\{x\},\qquad
 (X\setminus H)\times(H\cap C),\qquad \{y\}\times(H\cap B').
\]
Write $X=\{x_1,x_2,x_3,x_4\}$ and, when $\eta=5$,
$H\cap C=\{v_1,\dots,v_5\}$. The following requests, all of size at most
six, cover these pairs; together with $A,B,C,H$ there are at most seven
edges.

\emph{$H\cap X=\{x\}$.} Then $H\cap B'=\varnothing$. Let $x'\in X\setminus\{x\}$.
If $\eta\le1$, choose $S\subseteq C\setminus H$ with $|S|=2-\eta$ and use
$X\cup(H\cap C)\cup S$ and $\{x\}\cup(C\setminus(H\cup S))$. If $\eta=4$, use
$\{x,x'\}\cup(H\cap C)$, $(X\setminus\{x,x'\})\cup(H\cap C)$ and
$\{x\}\cup(C\setminus H)$. If $\eta=5$, use $X\cup\{v_1,v_2\}$,
$\{x\}\cup(C\setminus\{v_1,v_2\})$ and $(X\setminus\{x\})\cup\{v_3,v_4,v_5\}$.

\emph{$H\cap X=\varnothing$.} Put $Y_b=\{y\}\cup(H\cap B')$, of at most two
points. If $\eta\le1$, use $X\cup(H\cap C)$ and $Y_b$. If $\eta=4$, use
$\{x_1,x_2\}\cup(H\cap C)$, $\{x_3,x_4\}\cup(H\cap C)$ and $Y_b$. If
$\eta=5$, use $X\cup\{v_1,v_2\}$, $\{x_1\}\cup Y_b\cup\{v_3,v_4,v_5\}$ and
$\{x_2,x_3,x_4,v_3,v_4,v_5\}$.

In each case every candidate pair lies in one of the listed requests,
which completes the proof.
\end{proof}

\section{Three intersection gaps}\label{sec:gaps}

For the rest of the paper fix $k\ge7$ and write
\[
 h=\floor{k/7},\qquad s=k-7h,\qquad T=k-h=\ceil{6k/7},
\]
so that $h\ge1$, $0\le s\le6$ and $6k/7\le T\le k$. Assume
$\tau(\HH)>T$ for a finite $k$-uniform family $\HH$ with property
$(7,2)$. Define the integers
\begin{equation}\label{eq:endpoints}
 \alpha=\floor{T/2},\quad \beta=\floor{(4T-2k)/3},\quad
 \gamma=T-\ceil{k/2},\quad \delta=2k-T-2\beta,\quad \lambda=\delta-1.
\end{equation}
Normalized by $k$, these are approximately $3/7$, $10/21$, $5/14$,
$4/21$, $4/21$. The first stage excludes pair intersections in
$(\alpha,\min(\beta,\floor{k/2})]$. The other two stages, which exclude
$[\delta,\gamma]$ and then $(\beta,k/2]$, are needed only when
$\beta<\floor{k/2}$; only then are $\delta$ and $\lambda$ used, and only
then is $\lambda\ge0$.

\begin{lemma}[Endpoints]\label{lem:endpoints}
Put $r=\floor{(h+2s)/3}$. Then
\[
 \alpha=3h+\floor{s/2},\quad \beta=3h+r,\quad
 \gamma=\floor{(5h+s)/2},\quad \delta=2h+s-2r,\quad
 \lambda=2h+s-2r-1,
\]
and
\begin{equation}\label{eq:endpointfacts}
 0\le\gamma\le\alpha\le\beta,\qquad \beta\ge3h,\qquad
 \lambda\le\gamma-h,\qquad 4\lambda\le T,\qquad k-\beta-1+\lambda\le T .
\end{equation}
If $\beta<\floor{k/2}$, then $h<\delta\le\gamma$, so $0\le\lambda<\gamma<k$,
and moreover $T+\beta-k\le\gamma-1$.
\end{lemma}

\begin{proof}
The formulas follow from $k=7h+s$ and $T=6h+s$. Clearly $\gamma\ge0$, and
$\gamma\le\alpha$ follows from $T\le k$. As $T\ge4k/5$ we have
$(4T-2k)/3\ge T/2$, so $\beta\ge\alpha$; and $\beta\ge3h$ is clear. For
$\lambda\le\gamma-h=\floor{(3h+s)/2}$ it suffices that
$2r+1\ge\ceil{(h+s)/2}$; since $r\ge\floor{(h+s)/3}$, this follows from
$\ceil{n/2}\le2\floor{n/3}+1$ for integers $n\ge0$ (write $n=3a+b$ with
$0\le b\le2$; then $n\le4a+2$). Writing $h+2s=3r+e$ with $0\le e\le2$,
we get $4\lambda-T=2e-2r-s-4\le0$ and
$k-\beta-1+\lambda-T=s-3r-2\le-h-s\le0$.

If $\beta<\floor{k/2}$, then $2\beta<k$ gives $\delta>k-T=h$, and
$\delta=\lambda+1\le\gamma-h+1\le\gamma$. Finally
$T+\beta-k=\beta-h\le\floor{(7h+s)/2}-1-h=\floor{(5h+s)/2}-1=\gamma-1$.
\end{proof}

\begin{proposition}[First gap]\label{prop:gapone}
No two edges meet in $q$ points with $\alpha<q\le\min(\beta,\floor{k/2})$.
\end{proposition}

\begin{proof}
Apply Lemma~\ref{lem:caps} with cap $\gamma$: we have
$\gamma\le\floor{k/2}\le k-q$, and
$2k-T-2\gamma=3h+(k-2\floor{k/2})\le\alpha+1\le q\le2k-T$. The resulting
good triple has cells $x=q$ and $\gamma\ge y\ge z\ge0$. Put $d=y-z$,
$v=y+z$ and $S=x+v$. We use $2y\le2\gamma\le2T-k$, $3h<x\le\beta$, and
$\beta\le(4T-2k)/3\le3T-2k$ (the latter as $T\ge4k/5$).

\paragraph{Case 1: $d>x-h$.} Lemma~\ref{lem:splitA} applies to $(x,y,z)$.
From $v+d=2y\le2T-k$ we get $S\le x+2T-k-d<2T-k+h=T$. Next $2T\ge k+2y$.
Then $k+2x-y+z=k+2x-d<k+x+h\le2T$, as $x\le3T-2k$. Finally
$k+2x+y+3z=k+2x+2v-d<k+2x+4T-2k-3(x-h)=T+2k-x\le3T$, as $x\ge2h$.

\paragraph{Case 2: $d\le x-h$ and $v\le3T-k-2x$.} Lemma~\ref{lem:onecore}
applies with $(x,y,z)$ of the lemma equal to $(y,z,x)$. Its nine
conditions read
\begin{gather*}
 v\le2T-k\le T,\quad k+2y\le2T,\quad k+2z\le2T,\quad
 k+d-x\le k-h=T,\quad k-d-x\le T,\\
 3k-S\le3k-3h=3T,\quad 3k+S\le2k+x+2T\le5T,\quad k+v+2x\le3T,\quad
 2k+3x\le4T,
\end{gather*}
using $S\ge x\ge3h$, $x\ge h$ (for $k-d-x\le k-x$), $S\le x+2T-k$,
$x\le3T-2k$, the case hypothesis and $x\le\beta$.

\paragraph{Case 3: $v>3T-k-2x$.} Lemma~\ref{lem:symmetric} applies with
$m=x$, the largest cell. Indeed $k+x-v<2k+3x-3T\le T$; $k+x\le2T$ because
$x\le\beta\le2T-k$; $d=2y-v\le2T-k-v<2x-T$, so $2k+x+d<2k+3x-T\le3T$; and
$3k+S\le5T$ as in Case~2. Here $2k+3x-T\le3T$ is again $x\le\beta$.
\end{proof}

If $\beta\ge\floor{k/2}$, Proposition~\ref{prop:gapone} already excludes
every pair intersection in $(\alpha,k/2]$. For the next two stages we
therefore assume
\begin{equation}\label{eq:branch}
 \beta<\floor{k/2},
\end{equation}
so that Lemma~\ref{lem:endpoints} gives $h<\delta\le\gamma$ and
$0\le\lambda<\gamma$, and Proposition~\ref{prop:gapone} gives
$\Gap(\alpha,\beta)$.

\begin{proposition}[Middle gap]\label{prop:gaptwo}
Under~\eqref{eq:branch}, no two edges meet in $q$ points with
$\delta\le q\le\gamma$.
\end{proposition}

\begin{proof}
Apply Lemma~\ref{lem:caps} with cap $\beta$: $\beta<\floor{k/2}\le k-q$
and $0\le2k-T-q\le2k-T-\delta=2\beta$. Both new cells are at most $\beta$,
hence at most $\alpha$ by $\Gap(\alpha,\beta)$. So the good triple has
cells $\delta\le x\le\gamma$ and $\alpha\ge y\ge z\ge0$. This is the
structure of the finishing argument: $x$ is small and $y,z\le T/2$.

\paragraph{Case 1: $x+z\le k-\beta-1$.} Apply Corollary~\ref{cor:hall}
with $a=y$, $b=x$, $c=z$. Here $x\le\gamma$,
$z\le k-\beta-1-\delta=T+\beta-k-1\le\gamma$ by
Lemma~\ref{lem:endpoints}, and $x+z\le k-\beta-1\le k-3h-1<3T-2k$.

\paragraph{Case 2: $x+z\ge k-\beta$.} Apply Lemma~\ref{lem:twocores} with
$(x,y,z)$ of the lemma equal to $(y,x,z)$ and $\Gap(\alpha,\beta)$. Then
$g=(k-x-\beta)_+\le\min(y,z)$ because $z\ge k-\beta-x$ and $y\ge z$. The
first request has size $x+2g=\max(x,2k-x-2\beta)\le T$ because
$x\ge\delta$. The remaining conditions are $y+\alpha,\,z+\alpha\le2\alpha\le T$,
$k+2x\le k+2\gamma\le2T$, $3k\le4T$, and
$3k+S\le3k+\gamma+2\alpha\le\tfrac52k+2T\le5T$.
\end{proof}

\begin{proposition}[Upper gap]\label{prop:gapthree}
Under~\eqref{eq:branch}, no two edges meet in $q$ points with
$\beta<q\le\floor{k/2}$.
\end{proposition}

\begin{proof}
Apply Lemma~\ref{lem:caps} with cap $\gamma$, as in
Proposition~\ref{prop:gapone} (now $q\ge\beta+1\ge3h+1$). Both new cells
are at most $\gamma$, hence below $\delta$ by
Proposition~\ref{prop:gaptwo}; that is, $\Gap(\lambda,\gamma)$ holds and
the triple has cells
\[
 \beta<x\le\floor{k/2},\qquad \lambda\ge y\ge z\ge0 .
\]
Put $S=x+y+z$ and $J=2k-\gamma+2\floor{k/2}-3T=\floor{(h-s)/2}$. We use
$\lambda\le\gamma-h$, $4\lambda\le T$ and $k-x+y\le k-\beta-1+\lambda\le T$
from Lemma~\ref{lem:endpoints}.

\paragraph{Case 1: $S\le T$, and $S\le2T-k$ or $z<J$.}
Apply Lemma~\ref{lem:splitB} with $(x,y,z)$ of the lemma equal to
$(x,z,y)$. Its conditions are $S\le T$, $k+3x\le\tfrac52k\le3T$,
$k-x+y\le T$, and $k+2x+3z+y\le3T$. For the last, if $S\le2T-k$ then
$2x+3z+y\le2S\le4T-2k$ (as $y\ge z$), so $k+2x+3z+y\le4T-k\le3T$; if $z<J$ then $z<h/2$, $y\le\lambda\le2h+s$ and
$2x\le k$ give $k+2x+3z+y\le2k+\tfrac32h+2h+s\le18h+3s=3T$.

\paragraph{Case 2: $2T-k<S\le T$ and $z\ge J$.}
Apply Lemma~\ref{lem:gaptraces} with $\Gap(\lambda,\gamma)$. Here
\[
 S+p+t=\max(S,\ k-\gamma+z,\ k-\gamma+y,\ 2k-2\gamma-x)\le T,
\]
since $k-\gamma+y\le k-h=T$ and $2k-2\gamma-T=3h+(k-2\floor{k/2})\le\beta+1\le x$.
Similarly
\[
 k+3x+p+t=\max(k+3x,\ 2k-\gamma+2x-y,\ 2k-\gamma+2x-z,\ 3k-2\gamma+2x-S)\le3T:
\]
the first term is at most $\tfrac52k$; the middle two are at most
$2k-\gamma+2\floor{k/2}-J=3T$; and the last is at most
$3k-2\gamma+2\floor{k/2}-(2T-k)=6k-4T\le3T$, using
$\gamma=T-\ceil{k/2}$. Finally $y+z+2\lambda\le4\lambda\le T$.

\paragraph{Case 3: $S>T$.}
Apply Lemma~\ref{lem:gapcore} with $\Gap(\lambda,\gamma)$ and
$\Delta=S-T$. Write $\epsilon=k-2\floor{k/2}\in\{0,1\}$ and note
$\gamma-h=2T-k-\ceil{k/2}$. Using $y,z\le\lambda\le\gamma-h$ and
$x\le\floor{k/2}$,
\begin{align*}
 x+y&\le\floor{k/2}+\gamma\le T,\qquad k-x-z+\Delta=h+y\le\gamma,\\
 k-x-y&=h+z-\Delta<h+\lambda\le\gamma,\\
 2x+2\Delta+k-y-z&=4x+y+z+k-2T\le4\floor{k/2}+2(\gamma-h)+k-2T=2T-3\epsilon,\\
 y+z+2\lambda&\le4\lambda\le T,\\
 2k+x-y+\lambda+\Delta&=2k+2x+z+\lambda-T\le2k+2\floor{k/2}+2(\gamma-h)-T=3T-2\epsilon.
\end{align*}
\end{proof}

\section{Proof of the main theorem}\label{sec:proofmain}

\begin{proof}[Proof of Theorem~\ref{thm:main}]
For $2\le k\le6$ see the remark after Theorem~\ref{thm:main}. Let $k\ge7$
and suppose $\tau(\HH)>T=\ceil{6k/7}$. If $\beta\ge\floor{k/2}$,
Proposition~\ref{prop:gapone} shows that no pair intersection lies in
$(\alpha,k/2]$. Otherwise Propositions~\ref{prop:gapone}
and~\ref{prop:gapthree} show the same. Since $\alpha=\floor{T/2}$, every
pair intersection is at most $T/2$ or greater than $k/2$, and
Theorem~\ref{thm:finish} gives $\tau(\HH)\le T$, a contradiction.

Finally, let $k\ge2$. For families of nonempty sets of size at most $k$, add to each edge new
points private to that edge until it has $k$ points. A transversal of
the original family is one of the padded family. Conversely, replacing
each private point of a transversal of the padded family by an arbitrary
point of the corresponding original edge gives a transversal of the
original family of no larger size. So padding preserves the transversal
number, and it preserves property $(7,2)$ since a two-point transversal
of original edges meets the padded edges.
\end{proof}

\section{Where six-sevenths comes from}\label{sec:limits}

We record the three inequalities that determine the coefficient and
indicate why they do not appear to be relaxable by changing parameters.
Write $t=T/k$ and normalize sizes by $k$.

\emph{Static allocation.} Corollary~\ref{cor:hall} with $a=T/2$ needs
$3+t/2\le4t$, that is $t\ge6/7$; this is the load of the Hall
allocation, in which the distinguished cell lies in three of the four
requests and every private point in one.

\emph{Starting the first stage.} A request of $T$ points built from two
edges that meet in $q$ points leaves $2-t-q$ of their private points
available to the response. Both traces can be forced to be small (at
most $t-1/2$) only if $2-t-q\le2t-1$, that is $q\ge3-3t$. The first stage
must start where the finishing argument stops, at $t/2$, so
$3-3t\le t/2$, again $t\ge6/7$.

\emph{The two-medium triples.} In Proposition~\ref{prop:localfinish}, a
triple with two medium cells and a third cell $z\le1-t$ is closed by the
near-core construction, whose budget includes $2+M+z\le3t$ (in the
alternative $S>T$ of Lemma~\ref{lem:nearcore}); with $M=t/2$ and $z=1-t$
this is once more $t\ge6/7$. Cells with $z>1-t$ are closed by the
two-cores construction, whose first request needs $z\ge1-t$.

At $t=6/7$ the three thresholds coincide: $4t-3=t/2=3-3t=3/7$. For
$t<6/7$ they separate, and the window $4t-3<M<3-3t$ contains two kinds
of problematic configurations.
\begin{itemize}[leftmargin=2em]
\item If $M\le t/2$, the balanced request of Lemma~\ref{lem:balanced}
applied to a pair attaining the largest small intersection $M$ can
return triples with cells $M\ge y\ge z$, $t-\tfrac12<y\le1-\tfrac{t+M}2$
and $3t-2-M<z<1-t$; for $5/6<t<6/7$ these conditions are compatible
for every $M$ in the window with $M\le t/2$. For them both alternatives of Lemma~\ref{lem:nearcore} fail,
as does Lemma~\ref{lem:twocores}, in the orientations used in
Proposition~\ref{prop:localfinish}.
\item If $t/2<M<3-3t$, the first stage is not available, and any request
of $T$ points built from two edges meeting in $M$ points leaves
$2-t-M>2(t-\tfrac12)$ of their private points, so the response can
create a medium cell next to the cell of size $M>t/2$.
\end{itemize}
Exploratory computations, not used in any proof, are consistent with
this picture. They used continuous relaxations of the covering problems
(cell sizes normalized by $k$, request contents chosen by mixed integer
programming, adversarial responses bounded by column generation), at
$t=0.855$. Allowing oneself to discard one edge after the fourth
response and to close the resulting new good triple statically (the
device used in the small-maximum case of Theorem~\ref{thm:finish})
closed the tested configurations of the first kind. For configurations
of the second kind, the one-response constructions we tested, with or
without this device, required a budget above $0.855$ at every tested
point. Thus the obstruction moves but, as far as these experiments show,
reappears at the same coefficient. An improvement of $6/7$ seems to require a genuinely new
ingredient for triples that combine a cell slightly above $t/2$ with a
medium cell. The conjectured value $3/4$ remains open.

\appendix
\section{Proofs of the adaptive closures}\label{app:adaptive}

We use the notation of Section~\ref{sec:framework}. In each proof, $K$
is the response to the first request, and the three final requests are
described by their \emph{bases}, sets that they must contain; remaining
points are then distributed by Lemma~\ref{lem:allocation}. By
Lemma~\ref{lem:foursets} it suffices to cover the candidate pairs listed
there. When no triple cell survives, these are the three
\emph{complementary products}
\[
 (X\setminus K)\times(K\cap P_G),\qquad
 (Y\setminus K)\times(K\cap P_F),\qquad
 (Z\setminus K)\times(K\cap P_E).
\]

\begin{proof}[Proof of Lemma~\ref{lem:onecore}]
Request $K$ avoiding $X$, $Y$ and a set $Z_0\subseteq Z$ of
$\min(z,T-x-y)$ points; this is legal as $T\ge x+y$. Put
$Q=Z\cap K$, $A=(F\cap K)\setminus Q$, $B=(G\cap K)\setminus Q$,
$C=E\cap K$, $V=Z\setminus Q$ and $D=E\setminus(X\cup Y\cup C)$, with
sizes $q,a,b,c,v,d$. Then $A\subseteq P_F$, $B\subseteq P_G$,
$C\subseteq P_E$, $q\le(S-T)_+$, $q+a+b+c\le k$,
\[
 a\le k-x-z,\qquad b\le k-y-z,\qquad c\le k-x-y,\qquad d=k-x-y-c .
\]
By Lemma~\ref{lem:foursets} the candidate pairs are $Q\times E$,
$X\times B$, $Y\times A$ and $V\times C$. We make all three final
requests contain $Q$ and their union contain $E$, and cover the other
three products by bases.

Since $G\cap K=Q\cup B$ avoids $Y\cup Z_0$, we have $q+b\le k-y-|Z_0|$, so
$q+x+b\le\max(k+x-y-z,\ k+2x-T)\le T$; similarly
$q+y+a\le\max(k-x+y-z,\ k+2y-T)\le T$. Also $z\le4T/5\le T$, as
$4T\ge2k+3z\ge5z$.

\emph{Case 0: $c\le T-z$.} Use bases $Q\cup X\cup B$, $Q\cup Y\cup A$,
$Z\cup C$; each has at most $T$ points. Their total size plus $d$ is
$k+2q+a+b+z$. Using $a+b\le2k-x-y-2z$ and $q\le(S-T)_+$, this is at most
$\max(3k-S,\,3k+S-2T)\le3T$. Distribute $D$ among the residual
capacities (Lemma~\ref{lem:allocation}). All candidate pairs are
covered: $Q\subseteq Z$ lies in all three requests, and
$E=X\cup Y\cup C\cup D$ lies in their union.

In the remaining cases $c>T-z$. Put $a_0=k+x+z-2T$ and $b_0=k+y+z-2T$.

\emph{Case 1: $a<a_0$.} Split $C=C_2\sqcup C_3$ with $|C_2|\le T-y-a-z$ and
$|C_3|\le T-z$, and use bases
$Q\cup X\cup B$, $Y\cup A\cup Z\cup C_2$, $Z\cup C_3$. The split is
possible because $y+a+z<k+x+y+2z-2T\le T$ and
$a+c<(k+x+z-2T)+(k-x-y)$, so $c<(T-y-a-z)+(T-z)$ by $2k+3z\le4T$. The
total size plus $d$ is $k+q+a+b+2z\le2k+2z-c<2k+3z-T\le3T$; distribute
$D$. The products $X\times B$ and $Y\times A$ are covered by the first
two bases, $V\times C$ by the last two (both contain $Z\supseteq V$),
and $Q\times E$ as before.

\emph{Case 2: $b<b_0$.} Interchange the roles of $(X,B,x,b)$ and
$(Y,A,y,a)$ in Case~1.

\emph{Case 3: $a\ge a_0$ and $b\ge b_0$.} Put $u=c+z-T$, so $0<u\le c$.
Choose $C_{12}\subseteq C$ with $|C_{12}|=u$, put $C_3=C\setminus C_{12}$,
and choose an integer $\theta$ with
\[
 \max(0,\ y+a+z+u-T)\le\theta\le\min(v,\ T-q-x-b-u).
\]
Such $\theta$ exists: $q+b+c\le k-a\le2T-x-z$ gives $T-q-x-b-u\ge0$;
$q+a+c\le k-b\le2T-y-z$ gives $y+a+z+u-T\le v$; and
$y+a+z+u-T\le T-q-x-b-u$ is equivalent to $q+a+b+2c+x+y+3z\le4T$, which
follows from $q+a+b+c\le k$, $c\le k-x-y$ and $2k+3z\le4T$. Split
$V=V_{13}\sqcup V_{23}$ with $|V_{13}|=\theta$ and use bases
\[
 Q\cup X\cup B\cup C_{12}\cup V_{13},\qquad
 Q\cup Y\cup A\cup C_{12}\cup V_{23},\qquad Z\cup C_3 .
\]
The choice of $\theta$ bounds the first two by $T$ (the second has
$q+y+a+u+v-\theta=y+a+z+u-\theta$ points, as $q+v=z$), and the third has
exactly $T$ points. The total size plus $d$ is
$k+q+a+b+2z+u\le2k+3z-T\le3T$; distribute $D$. The first two bases cover
$X\times B$, $Y\times A$ and $V\times C_{12}$, the third covers
$V\times C_3$, and $Q\times E$ is covered as before.
\end{proof}

\begin{proof}[Proof of Lemma~\ref{lem:splitA}]
Request $K$ avoiding $X\cup Y\cup Z$ and $T-S$ points of $P_F$; this is
possible as $0\le T-S\le k-x-z$. No triple cell survives. Put
$A=F\cap K$, $B=G\cap K$, $C=E\cap K$, of sizes $a,b,c$; then
$a\le k-T+y$, $b\le k-y-z$ and $a+b+c\le k$. The candidate pairs are
$X\times B$, $Y\times A$, $Z\times C$.

If $a\le T-y-z$, split $B=B_1\sqcup B_2$ with $|B_i|\le T-x-z$ (possible
since $2(T-x-z)\ge k-y-z\ge b$) and use bases $X\cup Z\cup B_1$,
$X\cup Z\cup B_2$, $Y\cup Z\cup A$. Their total size plus $c$ is at most
$k+2x+y+3z\le3T$, so $C$ can be distributed. As every request contains
$Z$, all three products are covered.

If $a>T-y-z$, then $b+c\le k-a<k-T+y+z\le2(T-x-z)$. Split $B\cup C$ into
two parts of size at most $T-x-z$ and use $X\cup Z$ together with one
part, twice, and $Y\cup A$, of size at most $k-T+2y\le T$.
\end{proof}

\begin{proof}[Proof of Lemma~\ref{lem:splitB}]
Request $K$ avoiding $X\cup Y\cup Z$ and $p=\max(0,k-2(T-x)-y-z)$ points of
$P_G$. This fits ($p\le k-y-z$ as $T\ge x$), and the request has
$\max(S,k+3x-2T)\le T$ points. No triple cell survives; with
$A=F\cap K$, $B=G\cap K$, $C=E\cap K$ we have $b\le2(T-x)$,
$c\le k-x-y$ and $a+b+c\le k$.

If $b>k+y+z-T$, split $B$ into two parts of size at most $T-x$ and use
$X$ with each part, and $Y\cup Z\cup A\cup C$, of size at most
$y+z+k-b<T$.

Otherwise $b\le k+y+z-T\le2(T-x-y)$, the last step being
$3T\ge k+2x+3y+z$. Split $B=B_1\sqcup B_2$ with $|B_i|\le T-x-y$ and use
bases $X\cup Y\cup B_1$, $X\cup Y\cup B_2$, $Y\cup Z\cup C$; the last has
at most $k-x+z\le T$ points. Their total plus $a$ is at most
$k+2x+3y+z\le3T$; distribute $A$. Every request contains $Y$, which
covers $Y\times A$.
\end{proof}

\begin{proof}[Proof of Lemma~\ref{lem:gaptraces}]
Request $X\cup Y\cup Z$, $p$ points of $P_E$, $t$ points of $P_F$, and
$T-S-p-t$ points of $P_G$; these fit, as $u\ge0$ and $T\le k$. No
triple cell survives. The trace of $K$ on $E$ has at most
$k-x-y-p\le u$ points, hence at most $\ell$; likewise on $F$. Its trace
$B$ on $G$ has at most $k-T+x+p+t$ points. Split $B$ into two parts of
sizes at most $\ceil{(k-T+x+p+t)/2}$ and use bases $X\cup B_1$,
$X\cup B_2$ and $Y\cup Z\cup(F\cap K)\cup(E\cap K)$. The first two have
at most $T$ points because $k+3x+p+t\le3T$ (all quantities are
integers), and the third has at most $y+z+2\ell\le T$.
\end{proof}

\begin{proof}[Proof of Lemma~\ref{lem:gapcore}]
Request $K$ avoiding $X$, $Y$ and $\min(z,T-x-y)$ points of $Z$. Put
$P=Z\cap K$, $A=(F\cap K)\setminus P$, $B=(G\cap K)\setminus P$ and
$C=E\cap K$, with sizes $p,a,b,c$; then $p\le\Delta$ and $b\le k-y-z$.
Moreover $|E\cap K|=c\le k-x-y\le u$ and
$|F\cap K|=a+p\le k-x-z+\Delta\le u$, so $c\le\ell$ and $a+p\le\ell$ by the gap.
The candidate pairs are $P\times E$, $X\times B$, $Y\times A$ and
$(Z\setminus P)\times C$.

Split $B=B_1\sqcup B_2$ nearly equally and use bases $X\cup P\cup B_1$,
$X\cup P\cup B_2$ and $Y\cup Z\cup A\cup C$. The first two have at most
$x+\Delta+\ceil{(k-y-z)/2}\le T$ points, and the third at most
$y+z+2\ell-p\le T$. All three contain $P$. Distribute
$D=E\setminus(X\cup Y\cup C)$, so that the union contains $E$; the total
load is
\[
 k+x+z+2p+a+b\le k+x+z+p+\ell+b\le2k+x-y+\ell+\Delta\le3T .
\]
\end{proof}

\begin{proof}[Proof of Lemma~\ref{lem:twocores}]
Request $Y$, $g$ points of $X$, $g$ points of $Z$, and further points of
$F$ so that the request contains exactly $T-y$ points of $F$, taking
remaining points of $X\cup Z$ first. This is possible as
$2g\le T-y\le k$. Put $P=X\cap K$ and $R=Z\cap K$, of sizes $p$ and $r$;
by the priority rule $p+r\le(S-T)_+$. As $K$ avoids $Y$, these are the
only triple cells. Let $A=(F\cap K)\setminus(P\cup R)$,
$B=(G\cap K)\setminus R$, $C=(E\cap K)\setminus P$, with sizes $a,b,c$;
then $d:=|F\cap K|=p+r+a\le k-T+y$. The traces of $K$ on $E$ and on $G$
have at most $k-y-g\le u$ points, hence at most $\ell$.

Use bases $X\cup(G\cap K)$, $Z\cup(E\cap K)$ and $Y\cup(F\cap K)$, of
sizes at most $x+\ell$, $z+\ell$ and $k-T+2y$, all at most $T$. Each
contains $P$ and $R$. Distribute $D=P_E\setminus K$ and
$W=P_G\setminus K$, of sizes $k-x-y-c$ and $k-y-z-b$, so that the union
contains $E\cup G$. The total load is
\[
 2k-y+2(p+r)+a=2k-y+(p+r)+d\le3k-T+(S-T)_+\le3T,
\]
by $3k\le4T$ if $S\le T$ and by $3k+S\le5T$ otherwise. The candidate pairs
are $P\times G$ and $R\times E$, covered since $P,R$ lie in every request
and the union contains $E\cup G$, and $(X\setminus P)\times B$,
$(Z\setminus R)\times C$, $Y\times A$, covered by the three bases.
\end{proof}

\begin{proof}[Proof of Lemma~\ref{lem:nearcore}]
Request $K$ avoiding $Y\cup Z$ and $\min(x,T-y-z)$ points of $X$. The only
possible triple cell is $P=X\cap K$, of size $p\le\Delta$. Put
$X'=X\setminus P$, $A=K\cap P_E$, $B=K\cap P_F$, $C=K\cap G$ and
$W=P_G\setminus K$, with sizes $a,b,c,w$; then $C\subseteq P_G$,
\[
 a\le k-x-y,\qquad b\le k-x-z,\qquad c+w=k-y-z .
\]
By Lemma~\ref{lem:foursets} the candidate pairs are
$P\times(Y\cup Z\cup C\cup W)$, $X'\times C$, $Y\times B$ and $Z\times A$.

\emph{Case 1: $p+a\le m$.} Put $P$ in all three requests, $X'$ in
requests $1$ and $3$, $Y\cup B$ in request $1$, $Z\cup A$ in request $2$;
distribute $C$ between requests $1,3$ and $W$ among all three. The loads
before distribution are $x+y+b\le k+y-z\le T$, $p+a+z\le m+z\le T$ and
$x\le T$. The distribution of $C$ is possible since
$2x+y+b+c\le2k+x-2z\le2T$, and then that of $W$ since the total load is
$2x+k+p+a+b\le2k+x+m-z\le3T$; as $W$ may go to any request, these two
inequalities are all the conditions of Lemma~\ref{lem:allocation}. Every candidate pair is covered.

\emph{Case 2: $p+a>m$.} Then $|E\cap K|=p+a>k/2$ by the dichotomy (this
includes the possibility $K=E$), so $|G\cap K|=c<k/2$; in particular
$K\ne G$, and the dichotomy gives $c\le m$. Put $P$ in requests $1,2$; $X'$ and $C$ in
request $2$; $Y\cup B$ in request $1$; $Z$ in requests $1,3$; and
distribute $A$ between requests $1,3$ and $W$ between requests $1,2$. The
initial loads are $p+y+z+b\le\Delta+k-x+y\le T$, $x+c\le x+m\le T$ and
$z\le T$. The three conditions of Lemma~\ref{lem:allocation} for $A$ and
$W$ follow from
\begin{gather*}
 (p+y+z+b)+z+a\le\Delta+2k-2x+z\le2T,\\
 (p+y+z+b)+(x+c)+w=p+x+k+b\le\Delta+2k-z\le2T,\\
 p+y+2z+b+x+c+a+w=p+x+k+z+a+b\le\Delta+3k-x-y\le3T .
\end{gather*}
Every candidate pair is covered: $P$ meets $Y$ and $Z$ in request~$1$,
$C$ in request~$2$ and $W$ in request $1$ or $2$; $X'$ and $C$ share
request~$2$; $Y$ and $B$ share request~$1$; and both allowed labels of
$A$ meet the label $\{1,3\}$ of $Z$.

If $S\le T$, the first request contains $X$, so $P=\varnothing$ and
$\Delta=0$, and $W$ has no candidate partner. Case~1 then needs only the
first distribution, and Case~2 only the distribution of $A$, whose
conditions are $k-x+y\le T$ and $2k-2x+z\le2T$.
\end{proof}

\begin{proof}[Proof of Lemma~\ref{lem:twolarge}]
Put $Y=E\cap G$, $Z=F\cap G$, $Y'=F\cap H$, $Z'=E\cap H$ and
$U=Y\cup Z\cup Y'\cup Z'$; the six pair cells are disjoint. Write
$\sigma=|Y|+|Z|$ and $\rho=|Y'|+|Z'|$, both at most $2m$, and put
\[
 p=\ceil{k/2}-\min(\sigma,\rho),\qquad q=T-\ceil{k/2}-\max(\sigma,\rho).
\]
Both are nonnegative; $p\le\ceil{k/2}\le b$ and $q\le\floor{k/2}<x$.
Choose $B_0\subseteq B$ and $X_0\subseteq X$ of sizes $p,q$, and request
an edge $I$ avoiding $U\cup B_0\cup X_0$, a set of exactly $T$ points.
The trace of $I$ on $G$ has at most
$k-\sigma-p=\floor{k/2}+\min(\sigma,\rho)-\sigma\le\floor{k/2}$ points,
and similarly on $H$; by the dichotomy both are at most $m$. In
particular $|B\cap I|\le m$, and $|X\cap I|\le x-q$.

Choose $B_1\subseteq B$ with $B\cap I\subseteq B_1$ and
$|B_1|=\min(b,T-x)$, possible since $T-x\ge m$. Use the requests
$X\cup B_1$ and $(X\cap I)\cup(B\setminus B_1)$. The first has at most $T$
points; the second has $|X\cap I|+\max(0,b-T+x)$ points, and
$x-q+b-T+x\le T$ follows from $q\ge T-\ceil{k/2}-2m$ and the last
hypothesis.

As all triple intersections are empty, the candidate pairs of
$E,F,G,H$ lie in $X\times B$, $Y\times Y'$ or $Z\times Z'$. The last two
products lie in $U$, which $I$ avoids. A pair in $X\times B_1$ misses the
response to $X\cup B_1$. A pair $\{u,v\}$ with $u\in X$ and
$v\in B\setminus B_1$ has $v\notin I$, so it meets $I$ only if $u\in I$;
then it misses the response to $(X\cap I)\cup(B\setminus B_1)$. The seven
edges $E,F,G,H,I$ and the two responses form a bad subfamily.
\end{proof}

\section*{Research assistance}
This work made substantial use of AI systems. Claude (Anthropic) and
Codex (OpenAI) were used for mathematical exploration, for developing
and checking the local constructions and the case analysis, for
computational experiments that guided the choice of constructions, for
literature searches, and for preparing the manuscript. The proofs in
this paper are hand arguments with exact integer quantities; no
computation is needed to verify them. [The human authors must describe
here the independent verification they performed.]

\begin{thebibliography}{9}
\bibitem{Bloom}
T.~F.~Bloom, \emph{Erd\H{o}s problems}, Problem 644,
\url{https://www.erdosproblems.com/644}.

\bibitem{BKS}
M.~Buci\'c, D.~Kor\'andi and B.~Sudakov,
\emph{Covering graphs by monochromatic trees and Helly-type results for
hypergraphs}, Combinatorica \textbf{41} (2021), 319--352.

\bibitem{EFKT}
P.~Erd\H{o}s, D.~G.~Fon-Der-Flaass, A.~V.~Kostochka and Zs.~Tuza,
\emph{Small transversals in uniform hypergraphs},
Siberian Adv.\ Math.\ \textbf{2} (1992), 82--88.

\bibitem{EHT}
P.~Erd\H{o}s, A.~Hajnal and Zs.~Tuza,
\emph{Local constraints ensuring small representing sets},
J.\ Combin.\ Theory Ser.\ A \textbf{58} (1991), 78--84.

\bibitem{FKW}
D.~G.~Fon-Der-Flaass, A.~V.~Kostochka and D.~R.~Woodall,
\emph{Transversals in uniform hypergraphs with property $(7,2)$},
Discrete Math.\ \textbf{207} (1999), 277--284.

\bibitem{Kostochka}
A.~V.~Kostochka,
\emph{Transversals in uniform hypergraphs with property $(p,2)$},
Combinatorica \textbf{22} (2002), 275--285.
\end{thebibliography}
\end{document}
